\documentclass[11pt]{article}
\usepackage[T1]{fontenc}
\usepackage{lmodern}
\usepackage[margin=1in]{geometry}
\usepackage{amsmath,amssymb,amsthm}
\usepackage{microtype}
\usepackage{xurl}
\usepackage[colorlinks=true,linkcolor=blue,citecolor=blue,urlcolor=blue]{hyperref}
\hypersetup{pdftitle={A fixed birational map without a prime-field cycle-distribution limit},pdfauthor={Alec Kriebel}}
\newtheorem{theorem}{Theorem}
\newtheorem{lemma}[theorem]{Lemma}
\newtheorem{proposition}[theorem]{Proposition}
\theoremstyle{remark}
\newtheorem{remark}[theorem]{Remark}
\newcommand{\Fp}{\mathbf F_p}
\newcommand{\Q}{\mathbf Q}
\newcommand{\C}{\mathbf C}
\newcommand{\ord}{\operatorname{ord}}
\newcommand{\supp}{\operatorname{supp}}
\newcommand{\Li}{\operatorname{Li}}
\title{A fixed birational map without a\\prime-field cycle-distribution limit}
\author{Alec Kriebel\\[3pt]\small Independent researcher\\\small \href{https://orcid.org/0009-0001-9320-500X}{ORCID: 0009-0001-9320-500X}}
\date{Version 1.0 \quad 21 October 2026 (UTC)}
\begin{document}
\maketitle
\begin{abstract}
For the fixed birational map $L(u,v)=(u,(u^2+1)v)$, reduction at every
prime $p\equiv3\pmod4$ permutes all of $\Fp^2$. Its point-weighted
distribution of periods divided by $p$ has no limit. At every fixed
$1/2\le x<1$, the distribution function is
$1-\varphi(p-1)/(p-1)+o(1)$, with limsup $1$ and liminf at most $23/24$.
Moreover, neither these distribution-function values nor the period
measures converge on any further prime subset of positive upper relative
density. The latter statement uses the classical continuous empirical
law for $\varphi(p-1)/(p-1)$. The contribution is a precisely delimited
application to the broad existence assertion in the 2004 Oberwolfach
problem on finite-field maps. The triangular family, primitive-value
estimates, totient mean and scalar distribution law are established
ingredients; no new arithmetic theorem or historical firstness is claimed.
\end{abstract}

\section{Statement and scope}
Let
\begin{equation}\label{eq:map}
 L(u,v)=(u,(u^2+1)v),\qquad
 L^{-1}(u,v)=\left(u,\frac{v}{u^2+1}\right).
\end{equation}
These formulas define inverse rational maps over $\Q$ on a dense open
set. The forward map is polynomial, but its inverse is not polynomial
over $\C$. The rational invariant $I(u,v)=u$ has genus-zero generic
fibers, and $L^m(u,v)=(u,(u^2+1)^m v)$ shows that $L$ has infinite order.
The map belongs, after exchanging the coordinates, to the older
triangular polynomial family in Ostafe--Shparlinski
\cite[\S2.1, (1)--(4)]{OS}: take $m=1$, $a=1$, $b=0$,
$g_0(X_1)=X_1^2+1$, and $h_0=0$.

Set $\mathcal S=\{p\text{ prime}:p\equiv3\pmod4\}$. For $p\in\mathcal S$,
write $T_p(z)$ for the least positive period of an affine point and define
\begin{equation}\label{eq:measure}
 \mu_p=\frac1{p^2}\sum_{z\in\Fp^2}\delta_{T_p(z)/p},
 \qquad D_p(x)=\mu_p([0,x])
       =\frac1{p^2}\#\{z:T_p(z)\le px\}\quad(x\ge0).
\end{equation}
For a prime set $\mathcal A\subseteq\mathcal S$, positive upper relative
density in $\mathcal S$ means
\begin{equation}\label{eq:density}
 \limsup_{X\to\infty}
 \frac{\#\{p\le X:p\in\mathcal A\}}{\#\{p\le X:p\in\mathcal S\}}>0.
\end{equation}
Since $\mathcal S$ has density $1/2$ among primes, this is equivalent to
positive upper density relative to all primes.

\begin{theorem}\label{thm:main}
The map \eqref{eq:map} permutes all of $\Fp^2$ for every $p\in\mathcal S$.
Put $\theta_p=\varphi(p-1)/(p-1)$.
\par
\begin{enumerate}
\item For each fixed $x\in[1/2,1)$, eventually along $\mathcal S$,
\begin{equation}\label{eq:error}
 \left|D_p(x)-(1-\theta_p)\right|
       \le \frac1p+64p^{-1/4}.
\end{equation}
\item For every such fixed $x$,
\begin{equation}\label{eq:gap}
 \limsup_{\substack{p\to\infty\\p\in\mathcal S}}D_p(x)=1,
 \qquad
 \liminf_{\substack{p\to\infty\\p\in\mathcal S}}D_p(x)\le\frac{23}{24}.
\end{equation}
\item If $\mathcal A\subseteq\mathcal S$ satisfies \eqref{eq:density},
then $D_p(x)$ does not converge as $p\to\infty$ in $\mathcal A$, for any
fixed $x\in[1/2,1)$. The probability measures $\mu_p$ do not converge
weakly along $\mathcal A$ either.
\end{enumerate}
\end{theorem}

Vivaldi's 2004 contribution, joint work with Roberts
\cite[pp.2944--2945, Conjecture 1]{V}, uses the point weighting and scale
in \eqref{eq:measure}. It fixes the field-extension degree, permits
suitable positive-density reduction primes, and then takes $q=p$.
The stated existence assertion covers arbitrary birational maps;
genus one enters a subsequent heuristic. Here the fixed degree is one,
the forward polynomial acts on the full affine plane, and both directions
are defined at every point for $p\in\mathcal S$. Part (iii) also rules
out restoring the limit by a further positive-upper-density selection
inside this admissible class. Failure at $x=3/4$ avoids the reciprocal
integer thresholds in the proposed integrable step law. We draw no
conclusion about narrower elliptic-fiber or polynomial-automorphism
questions, or about the different prime-averaged statistic at scale $p^2$
in the source's later Conjecture 2.

\section{Exact fiber dynamics and primitive multipliers}
For $p\in\mathcal S$, $-1$ is a nonsquare, so $a_u=u^2+1\ne0$ for
every $u\in\Fp$. On each vertical fiber, $L$ multiplies the second
coordinate by $a_u$. This proves the permutation assertion. Every
$(u,0)$ has period one; every $(u,v)$ with $v\ne0$ has period
$\ord_{\Fp^\times}(a_u)$. In particular all points lie on cycles and all
periods divide $N=p-1$.

Define the number of primitive multipliers by
\[
 P_p=\#\{u\in\Fp:\ord_{\Fp^\times}(u^2+1)=N\}.
\]
Every proper divisor of $N$ is at most $N/2<p/2$. For a fixed
$x\in[1/2,1)$ and $p>1/(1-x)$, $N>px$. Thus exactly the primitive
fibers contribute excluded points, $p-1$ points per fiber:
\begin{equation}\label{eq:exact}
 D_p(x)=1-\frac{(p-1)P_p}{p^2}.
\end{equation}
This counts points, rather than assigning equal weight to cycles.
The threshold condition is essential: at $x=(p-1)/p$ all points are
included. Also $D_p(0)=0$ and $D_p(1)=1$. No estimate uniform for moving
$x\to1$ is asserted.

\begin{lemma}\label{lem:primitive}
For every $p\in\mathcal S$, with $\omega(N)$ the number of distinct
prime divisors of $N$,
\begin{equation}\label{eq:primitive}
 \left|\frac{P_p}{p}-\theta_p\right|
 \le\frac{\theta_p(2^{\omega(N)}-1)}{\sqrt p}.
\end{equation}
\end{lemma}
\begin{proof}
We give the familiar character argument to specify the input completely.
This mechanism is classical; a general quadratic primitive-value
formulation appears in Booker--Cohen--Sutherland--Trudgian
\cite[\S2, (5)--(6), Lemma 1]{BCST}.
Extend all multiplicative characters by zero at zero, including the
principal character. For $a\in\Fp^\times$, the primitive-element indicator is
\begin{equation}\label{eq:indicator}
 \boldsymbol1_{\ord(a)=N}
 =\theta_p\sum_{d\mid N}\frac{\mu(d)}{\varphi(d)}
             \sum_{\ord(\chi)=d}\chi(a).
\end{equation}
Indeed, write $a=g^k$ for a generator $g$. For a prime $\ell\mid N$,
the sum over characters of order $\ell$ is $\ell-1$ if $\ell\mid k$
and $-1$ otherwise. The squarefree-divisor sum factors into
$\prod_{\ell\mid N}(1-c_\ell(k)/(\ell-1))$. Multiplication by
$\theta_p$ gives zero if some $\ell\mid k$ and one otherwise, proving
\eqref{eq:indicator}.

Let $\eta$ be the quadratic character. For any nonprincipal $\chi$,
counting the square roots of $t$ and then substituting $t=-s$ gives
\begin{align}\label{eq:jacobi}
 S_\chi:=\sum_{u\in\Fp}\chi(u^2+1)
 &=\sum_{t\in\Fp}(1+\eta(t))\chi(t+1)\notag\\
 &=\eta(-1)J(\eta,\chi),\qquad
 J(\alpha,\beta)=\sum_{s\in\Fp}\alpha(s)\beta(1-s).
\end{align}
For completeness, for a nonprincipal character $\alpha$ its Gauss sum
$G(\alpha)=\sum_t\alpha(t)e^{2\pi i t/p}$ satisfies
\[
 |G(\alpha)|^2
 =\sum_{r\ne0}\alpha(r)\sum_{y\ne0}e^{2\pi i(r-1)y/p}
 =(p-1)-\sum_{r\ne0,1}\alpha(r)=p.
\]
If $\alpha,\beta,\alpha\beta$ are nonprincipal, decomposing the
product of Gauss sums by $t+s$ proves
$G(\alpha)G(\beta)=J(\alpha,\beta)G(\alpha\beta)$.
Hence $|J(\alpha,\beta)|=\sqrt p$. If $\beta=\alpha^{-1}$,
the substitution $w=s/(1-s)$ instead gives
$J(\alpha,\alpha^{-1})=-\alpha(-1)$, of modulus one.
Thus $|S_\chi|\le\sqrt p$, including the exceptional case $\chi=\eta$.

Sum \eqref{eq:indicator} over $a=u^2+1$. Its principal term is exactly
$p\theta_p$, since these values are nonzero. For each nontrivial
squarefree $d\mid N$, the $\varphi(d)$ characters cancel the denominator
in the triangle inequality. There are $2^{\omega(N)}-1$ such divisors,
which proves \eqref{eq:primitive}.
\end{proof}

There are six primes below $16$, and $2\le\ell^{1/4}$ for every prime
$\ell\ge17$. Consequently
\begin{equation}\label{eq:omega}
 2^{\omega(N)}\le64N^{1/4}.
\end{equation}
Combining \eqref{eq:exact}--\eqref{eq:omega}, and using
$0\le\theta_p\le1$, proves \eqref{eq:error}. In particular,
\begin{equation}\label{eq:bridge}
 D_p(x)=1-\theta_p+o(1)
 \quad\text{for every fixed }x\in[1/2,1),\quad p\to\infty\text{ in }\mathcal S.
\end{equation}
Alternatively the older estimate \cite[(5)--(6), Lemma 1]{BCST}
specializes directly with $Q(t)=t^2+1$, $e_1=1$, $e_2=N$.
Here $Q(0)=1$ is not primitive, so its $N(1,N)$ equals $P_p$;
the principal term is $p-1$ because its sum omits $t=0$.
That formulation gives the same $o(1)$ bridge. The stronger exact
principal term in the displayed elementary proof is not a new
character-sum theorem.

\section{Two explicit incompatible subsequences}
The arithmetic input for the second subsequence is the classical
unconditional mean
\begin{equation}\label{eq:mean}
 \frac1{\pi(X)}\sum_{p\le X}\theta_p\longrightarrow
 A:=\prod_{\ell\ \text{prime}}\left(1-\frac1{\ell(\ell-1)}\right).
\end{equation}
It is the $r=m=1$ specialization of the established shifted-prime
calculation in Menici--Pehlivan \cite[(1), Lemma 2, pp.262--263]{MP};
it also follows from Aivazidis--Sofos \cite[Theorem 1.1(i), $n=1$]{AS}.
It is not the primitive-root conjecture for a fixed base.

Here is a short proof of the needed mean, subject to the standard
Siegel--Walfisz theorem. M\"obius inversion gives
\[
 \sum_{p\le X}\theta_p
 =\sum_{d\le X-1}\frac{\mu(d)}d\,\pi(X;1,d).
\]
Set $K=(\log X)^6$. Since $\pi(X;1,d)\le X/d$, the tail $d>K$ is
$O(X/K)$. Uniformly for $d\le K$, Siegel--Walfisz gives
$\pi(X;1,d)=\Li(X)/\varphi(d)+O(X/(\log X)^4)$;
the summed error is $O(X\log K/(\log X)^4)=o(X/\log X)$.
The coefficient series $\sum_d1/(d\varphi(d))$ converges:
$\varphi(d)\ge\sqrt{d/2}$ follows by checking prime powers in
$\varphi(d)^2/d$ (only the factor $2^1$ can contribute below one).
The resulting absolutely convergent Euler product is $A$.
Finally $\pi(X)\sim\Li(X)$ proves \eqref{eq:mean}.
This is a restatement of a classical mean argument.

We need only a coarse positive bound. The factor at $2$ is $1/2$,
and
\[
 \sum_{\ell\ \text{odd prime}}\frac1{\ell(\ell-1)}
 \le\sum_{\substack{n\ge3\\n\ne4}}\frac1{n(n-1)}=\frac5{12}.
\]
The inequality $\prod(1-a_i)\ge1-\sum a_i$, first for finite products
and then by passage to the limit, yields
\begin{equation}\label{eq:Abound}
 A\ge\frac7{24}>\frac{13}{48}.
\end{equation}

Let $Q_j$ be the product of the first $j$ odd primes. The congruences
$p\equiv3\pmod4$ and $p\equiv1\pmod{Q_j}$ specify a reduced residue
class modulo $4Q_j$. By the Chinese remainder theorem and Dirichlet's
theorem, choose an increasing prime sequence $p_j$ in these classes.
Then
\[
 0\le\theta_{p_j}\le\frac12\prod_{\ell\mid Q_j}(1-1/\ell)
 \longrightarrow0
\]
by Euler's vanishing prime product. Thus \eqref{eq:bridge} gives
$D_{p_j}(x)\to1$ for every fixed $x\in[1/2,1)$.

For every odd prime, $\theta_p\le1/2$. If only finitely many
$p\in\mathcal S$ satisfied $\theta_p\ge1/24$, the prime number theorem
in the two odd classes modulo $4$ would give
\[
 \limsup_{X\to\infty}\frac1{\pi(X)}\sum_{p\le X}\theta_p
 \le\frac12\cdot\frac12+\frac12\cdot\frac1{24}
 =\frac{13}{48},
\]
contradicting \eqref{eq:mean}--\eqref{eq:Abound}. An increasing
sequence $q_j\in\mathcal S$ therefore satisfies
$\theta_{q_j}\ge1/24$. Equation \eqref{eq:bridge} gives
$\limsup_jD_{q_j}(x)\le23/24$, proving \eqref{eq:gap}.
The two sequences are independent of $x$; the exact identity is
eventual for each fixed $x$. The numerical bound is deliberately coarse
and is not offered as a new arithmetic refinement.

\section{Positive-density selections and weak convergence}
The two subsequences alone do not exclude convergence on every further
positive-density selection. For that quantifier we use the following
older scalar distribution theorem.

\medskip\noindent\textbf{Classical scalar input.}
There is an atomless probability measure $\nu$ on $[0,1]$ such that
\begin{equation}\label{eq:scalar}
 \frac1{\pi(X)}\sum_{p\le X}\delta_{\theta_p}
 \ \Longrightarrow\ \nu\qquad(X\to\infty).
\end{equation}
Equivalently, the empirical distribution function of $\theta_p$ is
continuous in the limit. This is the shifted-prime law associated
with K\'atai. It is recorded directly for this scalar in
Deshouillers--Hassani \cite[\S\S1--2, Proposition 2.1 with $m=2$]{DH},
and in Aivazidis--Sofos \cite[Theorem 1.3(i), $n=1$, and \S6]{AS}.
In the latter formulation the theorem is stated over prime powers;
proper powers number $O(\sqrt X\log X)=o(\pi(X))$, so the same
empirical law holds over primes, as its proof explicitly explains.
Continuity, rather than absolute continuity, is the property used here.
No hypothesis such as GRH is imposed on this input.

\begin{lemma}\label{lem:atom}
If a bounded real sequence $b_p$, indexed by primes, has an atomless
limiting empirical law, it cannot converge along a prime subset of
positive upper density relative to the primes.
\end{lemma}
\begin{proof}
Suppose $b_p\to c$ on a set $\mathcal A$ with upper density $\delta>0$.
For each $\varepsilon>0$, all but finitely many members of
$\mathcal A$ have $b_p\in[c-\varepsilon,c+\varepsilon]$.
Choose endpoints $X_j\to\infty$ realizing the upper density.
The limiting empirical measure assigns the closed interval at least
$\delta$: its endpoints have measure zero, so its frequencies converge.
Letting $\varepsilon\downarrow0$ and using continuity from above forces
an atom of mass at least $\delta$ at $c$, a contradiction.
\end{proof}

If $D_p(x)$ converged along a set satisfying \eqref{eq:density}, then
\eqref{eq:bridge} would force $\theta_p$ to converge there. The density
of $\mathcal S$ and Lemma~\ref{lem:atom}, applied to \eqref{eq:scalar},
exclude that possibility.

To treat weak convergence without a choice of continuity points, use
the bounded continuous function on $[0,\infty)$
\[
 h(t)=\max\{0,\min\{1,2t-1\}\}.
\]
For a proper-divisor period $d\mid p-1$, $d/p<1/2$ and $h(d/p)=0$.
A maximal period contributes $h((p-1)/p)=1-2/p$. Consequently the
exact identity
\begin{equation}\label{eq:moment}
 \int h\,d\mu_p
 =\left(1-\frac2p\right)\frac{(p-1)P_p}{p^2}
 =\theta_p+o(1)
\end{equation}
holds as $p\to\infty$ in $\mathcal S$, by
Lemma~\ref{lem:primitive} and \eqref{eq:omega}.
Weak convergence along a positive-upper-density subset would force
the left side to converge. Lemma~\ref{lem:atom} again gives a
contradiction. This proves part (iii) of Theorem~\ref{thm:main}.
In particular, both conclusions survive deletion of any set of prime
density zero. No assertion about arbitrary zero-density selections
is required.

\section{Attribution, historical boundary, and reproducibility}
The mechanism is a short application of established results. The
specific assertion documented here combines the exact point-period
identity \eqref{eq:exact} with classical arithmetic inputs for this
fixed map and the statistic in \eqref{eq:measure}. A bounded search
of the closest inspected finite-field dynamics sources did not
authenticate an earlier statement of this exact application. This
does not certify worldwide absence, historical firstness, or present
open-problem status.

A specific unresolved bibliographic lead must be retained.
Roberts \cite[p.220, reference 9]{R} cites N.~Neum\"arker,
J.~A.~G.~Roberts, C.-M.~Viallet and F.~Vivaldi,
\emph{The dynamics of reversible birational maps over finite fields},
UNSW Preprint (2011). The body and its version relationships were
not recovered in the bounded audit. Citation contexts cannot
determine whether that unseen manuscript contains the present
application. If an exact prior statement is later authenticated,
this note's historical description requires reconciliation; none
of the mathematical deductions depend on assuming its absence.

The accompanying source and verification archive contains the TeX
source, a dependency-free exact finite checker, source pointers,
and byte-preserved original research files with manifests. The
new checker passes $147{,}619$ exact finite checks over the $87$
admissible primes through $1000$, including eight complete affine orbit
graphs with $6{,}088$ points and six bad-reduction controls. It verifies
fiber orders, divisor fixed-point counts, the
point-weighted CDF identity and the continuous moment, and tests
the endpoint and bad-reduction controls. Its run receipt specifies
the actual range and assertion count. Historical AI audit receipts
also record $8{,}347{,}742$ finite assertions and $230{,}780$ analytic
controls; they are preserved as earlier runs, not attributed to the
new checker. The original SymPy-dependent scripts are retained;
this preparation does not claim a successful replay of those scripts
in an environment without SymPy. Finite computations check finite
identities and implementation consistency. The infinite-prime
conclusions follow from the proofs and the explicitly cited
classical inputs, not from numerical extrapolation.

\medskip\noindent\textbf{AI disclosure and review status.}
AI systems were extensively involved in the research, mathematical
model and proof development, computation, literature investigation,
manuscript writing, and adversarial verification. The author is
responsible for the submitted claims. This is an unrefereed preprint;
it has not undergone formal independent human peer review.
Internal AI reviews are not human peer review.

\begin{thebibliography}{9}
\bibitem{V} F.~Vivaldi (joint work with J.~A.~G.~Roberts),
\emph{Maps over finite fields: integrability and reversibility},
Oberwolfach Reports 54/2004, pp.2944--2945.
\url{https://webspace.maths.qmul.ac.uk/f.vivaldi/research/Oberwolfach.pdf}.

\bibitem{OS} A.~Ostafe and I.~E.~Shparlinski,
\emph{On the degree growth in some polynomial dynamical systems and
nonlinear pseudorandom number generators}, arXiv:0902.3884v3 (2009).
\url{https://arxiv.org/abs/0902.3884v3}.

\bibitem{BCST} A.~R.~Booker, S.~D.~Cohen, N.~Sutherland and T.~Trudgian,
\emph{Primitive values of quadratic polynomials in a finite field},
arXiv:1803.01435v2 (2018).
\url{https://arxiv.org/abs/1803.01435v2}.

\bibitem{MP} L.~Menici and C.~Pehlivan,
\emph{Average $r$-rank Artin conjecture}, Acta Arithmetica
\textbf{174}(3) (2016), 255--276.
\href{https://doi.org/10.4064/aa8258-4-2016}{doi:10.4064/aa8258-4-2016}.

\bibitem{AS} S.~Aivazidis and E.~Sofos,
\emph{On the distribution of the density of maximal order elements in
general linear groups}, arXiv:1308.5701v2 (2014).
\url{https://arxiv.org/abs/1308.5701v2}.

\bibitem{DH} J.-M.~Deshouillers and M.~Hassani,
\emph{A note on the distribution function of $\varphi(p-1)/(p-1)$},
Journal of the Australian Mathematical Society \textbf{93} (2012),
77--83; published online 16 January 2013.
\href{https://doi.org/10.1017/S1446788712000250}{doi:10.1017/S1446788712000250}.

\bibitem{R} J.~A.~G.~Roberts,
\emph{Order and symmetry in birational difference equations and their
signatures over finite phase spaces}, Proceedings of the International
Workshop Future Directions in Difference Equations, Vigo (2011),
213--221, especially reference 9.
\url{https://www.dma.uvigo.es/~eliz/pdf/Roberts.pdf}.
\end{thebibliography}
\end{document}
